In this problem, we will explore the connection between the physical properties
of exoplanets and their host stars and will use the observational data to
discover as much as possible about these systems. You may neglect interstellar
extinction.

\medskip
\noindent\textbf{Part 1 (20 points).}

\begin{center}
\begin{tabular}{cccccc}
\toprule
Name of & Name of & $T_{\mathrm{eff}}$ & $g$ & $m_v$ & parallax \\
planet & star & (K) & (m\,s$^{-2}$) & (magnitudes) & (milliarcsec) \\
\midrule
Gorgona & HD\,209458 & 5980 & 347 & 7.63 & 20.67 \\
\bottomrule
\end{tabular}

\smallskip
Table 1: Observational data for exoplanet Gorgona and its parent star HD\,209458
\end{center}

The effective temperature ($T_{\mathrm{eff}}$) and the gravitational
acceleration in the surface of the star ($g$) can be measured from the shape of
the spectrum and its absorption lines.
% sic: "in the surface" as in the original
The visual apparent magnitude ($m_v$) and parallax are measured by doing
photometry and astrometry, respectively.

Additionally, it has been observed that every $3.52$ days the brightness of the
star drops due to the transit of the planet in front of it, as it is
represented in this lightcurve:

\ioaafig{2021_DA_DA2-f1.pdf}
% f1 = Figure 1: Normalized Flux (y-axis) against time (x-axis) for the parent
% star HD 209458

\begin{description}
\item[(1.1)] Use the given information to calculate the following quantities
  for the HD\,209458 system: \hfill[20\,pt]

  \begin{center}
  \begin{tabular}{ccccc}
  \toprule
  Luminosity of & Radius of & Mass of & Mean planet's & Radius of the planet\\
  the star & the star & the star & orbital radius & in Jupiter's radius\\
  $L_\star$ $[L_\odot]$ & $R_\star$ $[R_\odot]$ & $M_\star$ $[M_\odot]$
    & $a$ $[\mathrm{au}]$ & $R_p$ $[R_J]$\\
  \midrule
  \phantom{xxxx} & \phantom{xxxx} & \phantom{xxxx} & \phantom{xxxx}
    & \phantom{xxxx}\\
  \bottomrule
  \end{tabular}
  \end{center}

  \textbf{Note:} Assume that the bolometric correction for all F and G type
  stars is the same.
\end{description}

\medskip
\noindent\textbf{Part 2 (25 points).}

The habitable zone is defined as the region in which a planet may have liquid
water on its surface. This is mainly related to the amount of radiation
received from the host star, which must be within some limits to ensure a
favorable range of planet surface temperatures.

We define the effective flux received by a planet as
$S_{\mathrm{eff}} = \frac{L}{a^2}$, where $L$ is the star luminosity in solar
units, and $a$ is the mean orbital radius in au. The minimum flux in the
habitable zone can be approximated by
\[
S_{\mathrm{min}} = S_{\mathrm{eff}_\odot} + n\cdot T_\star + b\cdot T_\star^2
  + c\cdot T_\star^3 + d\cdot T_\star^4,
\]
where $T_\star = (T_{\mathrm{eff}} - T_{\mathrm{eff}_\odot})$, and
$S_{\mathrm{eff}_\odot}$ is the equivalent flux for the case of the Sun, which
along with the coefficients $n$, $b$, $c$, $d$ is given in the following
table. The maximum flux for habitability, $S_{\mathrm{max}}$, is found with the
same equation but different constants:

\begin{center}
\begin{tabular}{ccc}
\toprule
Constant & $S_{\mathrm{max}}$ & $S_{\mathrm{min}}$\\
\midrule
$S_{\mathrm{eff}_\odot}$ & $1.0512$ & $0.3438$\\
$n$ & $1.3242\times10^{-4}$ & $5.8942\times10^{-5}$\\
$b$ & $1.5418\times10^{-8}$ & $1.6558\times10^{-9}$\\
$c$ & $-7.9895\times10^{-12}$ & $-3.0045\times10^{-12}$\\
$d$ & $-1.8328\times10^{-15}$ & $-5.2983\times10^{-16}$\\
\bottomrule
\end{tabular}
\end{center}

The table below gives real data for 7 different star-planet systems. However
planet names have been changed to honour some natural sanctuaries in Colombia:

\begin{center}
\begin{tabular}{cc cc}
\toprule
\multicolumn{2}{c}{Stellar Parameters} & \multicolumn{2}{c}{Planet Parameters}\\
$T_{\mathrm{eff}}$ [K] & $M_V$ [mag] & Name & $a$ [au]\\
\midrule
6180 & 3.68 & Tayrona   & 0.04\\
5730 & 3.87 & Iguaque   & 0.04\\
5980 & 4.21 & Gorgona   & 0.04\\
5480 & 6.04 & Amacayacu & 0.08\\
5770 & 3.48 & Malpelo   & 0.05\\
6130 & 3.07 & Pisba     & 0.03\\
6140 & 3.85 & Tatam\'a  & 0.06\\
\bottomrule
\end{tabular}
\end{center}
% \ioaafig{2021_DA_DA2-f2.pdf} % f2 = image of the 7-system table transcribed
% just above; kept here (commented) to avoid printing the same table twice

\begin{description}
\item[(2.1)] In the following figure, the vertical axis represents the
  effective temperature of stars, and the horizontal axis represents the
  effective flux received by orbiting planets. The dot marked on the graph
  represents planet Earth, and dashed lines mark the limits of the habitable
  zone. \hfill[15\,pt]

  % NOTE: the original shows here an unlabelled grid (T_eff versus S_eff) with
  % the Earth marked as a dot and the habitable-zone limits as dashed lines;
  % no figure was extracted for it in the crop set (f1-f5).

  Put numerical labels at tickmark position on both axes. Draw on the same
  figure the exact position where Gorgona and Amacayacu would be, as if they
  were also located at 1\,au from their corresponding stars.
\item[(2.2)] Now considering the real orbital radius given in the table for
  each planet, indicate with YES or NO which of them are in the habitable
  zone. Show your quantitative reasoning on the working
  sheets. \hfill[10\,pt]

  \begin{center}
  \begin{tabular}{cc}
  \toprule
  Planet's Name & In habitable zone? YES / NO\\
  \midrule
  Tayrona   & \\
  Iguaque   & \\
  Gorgona   & \\
  Amacayacu & \\
  Malpelo   & \\
  Pisba     & \\
  Tatam\'a  & \\
  \bottomrule
  \end{tabular}
  \end{center}
  % \ioaafig{2021_DA_DA2-f3.pdf} % f3 = image of the blank YES/NO answer table
  % transcribed just above; kept here (commented) to avoid duplication
\end{description}

\medskip
\noindent\textbf{Part 3 (30 points).}

In the last page you find a list of 38 exoplanets, and the goal is to find out
if low-mass exoplanets (LME) and high-mass exoplanets (HME) tend to orbit
around stars with different characteristics.

\begin{description}
\item[(3.1)] To get a robust low-mass subsample one can apply a technique
  called ``iterative sigma-clipping''. The idea is to compute the mean ($\mu$)
  and the standard deviation ($\sigma$) of the masses and to exclude from the
  sample, those planets with masses above $\mu + \sigma$. Then repeat the same
  steps with the remaining subsample two more times. We will say that planets
  in the final subsample are the low-mass ones, and those excluded during the
  iterations, the high-mass ones. Fill the following table with the numbers
  you find in the process: \hfill[10\,pt]

  \begin{center}
  \begin{tabular}{cccccc}
  \toprule
  Low-mass sample & Sample size & $\mu$ & $\sigma$ & $\mu+\sigma$
    & No.\ of planets\\
  & & & & & to exclude\\
  \midrule
  Full / Original & 38 & & & & \\
  subsample after 1st iteration & & & & & \\
  subsample after 2nd iteration & & & & & \\
  subsample after final iteration & & --- & --- & --- & ---\\
  \bottomrule
  \end{tabular}
  \end{center}
  % \ioaafig{2021_DA_DA2-f4.pdf} % f4 = image of the blank sigma-clipping
  % answer table transcribed just above; kept here (commented) to avoid
  % duplication
\item[(3.2)] Make a plot using the X-axis for the serial number of the planets
  in the list $(1, 2, 3, \ldots)$, and the Y-axis for the mass of the
  planets. Mark 3 horizontal lines at the $\mu + \sigma$ thresholds you found
  in the iterations. \hfill[5\,pt]
\item[(3.3)] Let's investigate the possible difference in the effective
  temperatures of host stars in both groups, computing some descriptive
  statistics: \hfill[10\,pt]

  \begin{center}
  \begin{tabular}{cccccc}
  \toprule
  & Min. & 1st Quartile & Median (2nd quartile) & 3rd Quartile & Max\\
  \midrule
  LME & & & & & \\
  HME & & & & & \\
  \bottomrule
  \end{tabular}
  \end{center}
\item[(3.4)] Draw boxplots summarizing the numbers you just computed. Do you
  see a clear difference in the temperatures of the stars orbited by low-mass
  and high-mass planets? Write YES or NO. \hfill[5\,pt]

  % NOTE: the original provides here an empty chart frame labelled
  % "Planetary type vs host star temperature (K)", with HME and LME on the
  % vertical axis and temperature ticks 5200-6400 K on the horizontal axis;
  % no figure was extracted for it in the crop set (f1-f5).
\end{description}

\medskip
\noindent The list of 38 exoplanets (planet mass in Jupiter masses,
$T_{\mathrm{eff}}$ of the host star in K):

\begin{center}
\begin{tabular}{clcc}
\toprule
No. & Name of Planet & Planet Mass $[M_J]$ & $T_{\mathrm{eff}}$ of the Host Star\\
\midrule
1 & KEPLER-37 b & 0.01 & 5520\\
2 & KEPLER-21 b & 0.02 & 6256\\
3 & HD 97658 b  & 0.02 & 5468\\
4 & HD 46375 b  & 0.23 & 5345\\
5 & HD 219134 h & 0.28 & 5209\\
6 & HD 88133 b  & 0.30 & 5582\\
7 & HD33283 b   & 0.33 & 5877\\
8 & HD 149026 b & 0.36 & 6096\\
9 & BD-10 3166 b & 0.46 & 5578\\
10 & HD 75289 b & 0.47 & 6196\\
\bottomrule
\end{tabular}
\end{center}
% TABLE ABRIDGED — full data in crop edition
% (remaining rows: 11 HD 217014 b 0.47 5755; 12 HD 2638 b 0.48 5564;
% 13 WASP-13 b 0.49 6025; 14 WASP-34 b 0.59 5771; 15 HD 209458 b 0.69 5988;
% 16 HAT-P-30 b 0.71 6177; 17 WASP-76 b 0.92 6133; 18 WASP-74 b 0.97 5727;
% 19 HAT-P-6 b 1.06 6442; 20 HD189733 b 1.14 5374; 21 WASP-82 b 1.24 6257;
% 22 KELT-7 b 1.29 6460; 23 HD 149143 b 1.33 6067; 24 KELT-3 b 1.42 6404;
% 25 KELT-2A b 1.49 6164; 26 HD86081 b 1.50 6015; 27 HAT-P-7 b 1.74 6270;
% 28 HD 118203 b 2.14 5847; 29 HAT-P-14 b 2.20 6490; 30 WASP-38 b 2.71 6178;
% 31 HD17156 b 3.20 5985; 32 KELT-6 c 3.71 6176; 33 HD 75732 d 3.86 5548;
% 34 HD 115383 b 4.00 5891; 35 HD 120136 b 5.84 6210; 36 WASP-14 b 7.34 6195;
% 37 HAT-P-2 b 8.74 6439; 38 XO-3 b 11.79 6281)

\ioaafig{2021_DA_DA2-f5.pdf}
% f5 = image of the full 38-planet list (all rows), kept as the complete
% data table
